\section{Run details}
\label{app:runs}

\subsection{Five-window paired training deltas}

Table~\ref{tab:windows} partitions each run into five 536{,}870{,}912-token
windows (1{,}024 updates each) and reports the mean per-update paired
difference against the seed-matched baseline in each. This answers a different
question from the end-of-run means in Table~\ref{tab:main}: whether an arm's
advantage is a stable offset or a transient.

The shapes separate into three kinds. Arms with dense left factors open with a
large advantage that decays by a factor of three --- \dense{} runs from
$-0.015940$ to $-0.004975$ --- which is consistent with a benefit that matters
most while curvature is changing fastest. Arms with diagonal left factors are
nearly flat, and \alldiag{} is actually strongest in the middle of training.
The matched nulls begin \emph{worse} than the baseline and end slightly better,
which is the signature of an accumulating execution-path difference rather than
a geometric one. Arms that degrade the input side diverge monotonically and
never recover, with \texttt{inner\_outer\_all\_diag} at $+0.058$ in the first
window.

\begin{table}[t]
\centering
\caption{Mean per-update paired training-loss difference against the
seed-matched baseline, in five consecutive 536{,}870{,}912-token windows.
Negative is better. Windows are exactly comparable across arms within a seed
because the token order is pinned. Generated by \texttt{scripts/analyze.py}.}
\label{tab:windows}
\footnotesize\setlength{\tabcolsep}{3.4pt}
% Generated by scripts/analyze.py. Do not edit.
\begin{tabular}{lrrrrr}
\toprule
Arm & Window 1 & Window 2 & Window 3 & Window 4 & Window 5 \\
\midrule
mlpdiag & $-0.0096$ & $-0.0066$ & $-0.0057$ & $-0.0051$ & $-0.0051$ \\
dense, both widths & $-0.0159$ & $-0.0078$ & $-0.0059$ & $-0.0049$ & $-0.0050$ \\
dense, lagged & $-0.0178$ & $-0.0078$ & $-0.0058$ & $-0.0054$ & $-0.0057$ \\
all diagonal & $-0.0038$ & $-0.0047$ & $-0.0054$ & $-0.0042$ & $-0.0037$ \\
outer factor only & $-0.0034$ & $-0.0045$ & $-0.0046$ & $-0.0041$ & $-0.0039$ \\
no 2048 output & $-0.0002$ & $-0.0025$ & $-0.0023$ & $-0.0021$ & $-0.0021$ \\
down only & $-0.0006$ & $+0.0009$ & $+0.0000$ & $-0.0007$ & $-0.0012$ \\
\addlinespace
\textbf{null (dense)} & $+0.0017$ & $-0.0005$ & $-0.0021$ & $-0.0013$ & $-0.0014$ \\
\textbf{null (mlpdiag)} & $+0.0012$ & $+0.0001$ & $-0.0005$ & $-0.0006$ & $-0.0007$ \\
\addlinespace
mlpdiag, seed 1 & $-0.0095$ & $-0.0050$ & $-0.0041$ & $-0.0043$ & $-0.0035$ \\
mlpdiag, seed 2 & $-0.0081$ & $-0.0048$ & $-0.0038$ & $-0.0032$ & $-0.0035$ \\
\addlinespace
Muon & $+0.0930$ & $+0.0162$ & $+0.0063$ & $+0.0030$ & $+0.0062$ \\
output only & $+0.0301$ & $+0.0219$ & $+0.0159$ & $+0.0121$ & $+0.0079$ \\
both sides diagonal & $+0.0585$ & $+0.0336$ & $+0.0286$ & $+0.0224$ & $+0.0134$ \\
\bottomrule
\end{tabular}

\end{table}

\subsection{The fixed-age and systems waves}
\label{app:systems}

Four later waves changed implementation rather than geometry, and we report
them separately because none has a matched null-output arm, so none isolates
geometry from execution. They are nevertheless the only place where the
diagonal-output geometry was measured under a controlled host-refresh policy.

\paragraph{Fixed-age input refresh.}
One wave pinned input host-basis installation to the next update boundary, so
that every logged rank-0 installation has age exactly one and no pending
refresh is discarded --- 23{,}712 to 23{,}765 installed events per arm and zero
discards. Under that policy, against its own same-wave baseline, the
all-diagonal, outer-diagonal and fused-outer arms improve final C4 by
$0.003831$, $0.002679$ and $0.002188$ nats at $+9.98$\%, $+9.77$\% and
$+8.86$\% main-function time. This removes variable \emph{input} staleness as an
explanation for the output-side benefit. It says nothing about output staleness,
because these arms' diagonal output factors were already recomputed every
update by construction.

\paragraph{Fused and combined products.}
Fusing the post-Newton--Schulz products from four right multiplications to
three saves 30{,}064{,}771{,}072 multiply--add operations per rank per update in
theory, and 5.4--15.5\% in isolated kernel benchmarks. At whole-model scale it
saved nothing measurable: the combined arm ran 9.87 seconds \emph{slower} than
the unfused one over a 6{,}470-second run, improved final C4 by $0.000856$ and
worsened final selection by $0.001216$. We include this because it is a clean
example of a kernel-level saving that does not survive to end-to-end time, and
because it is the reason Section~\ref{sec:program} insists on launch-anchored
timing rather than product counts.

\paragraph{Upstream-product variant and its null.}
A systems wave compared the project's own right-product implementation against
the pinned upstream one, both with fused diagonal output factors, plus a matched
null using upstream arithmetic. The upstream variant saved $0.705$\% of main
time against the legacy one and still cost $7.64$\% over the baseline. Its
margin over its own matched null was $0.000383$ nats of final selection loss and
$0.000443$ of late selection --- the smallest attributable margin anywhere in
this study, and a direct illustration of why the null matters: the null alone
accounts for $-0.001814$ final C4 and $-0.001404$ final selection against the
baseline.

\paragraph{Selected-rate transfer.}
The wave discussed in Section~\ref{sec:tuning} carried screening-selected
learning rates to the full budget. Its arms are the only ones in this study not
at the recipe rate, and its outcome ($+0.049$ and $+0.050$ selection for two
geometries, $-0.0029$ for the third) is reported there rather than here because
it bears on the paper's central claim.

\subsection{Waves that produced no optimizer result}

Three waves are recorded as failures and contribute no numbers to any table. We
list them because their failure modes are informative about the difficulty of
this kind of study, and because omitting them would misrepresent the ratio of
attempted to completed work.

\begin{itemize}
\item A fused-diagonal candidate failed an A100 direction-parity gate before
  training: 606 of 1551 entries on one rank differed from the reference by up to
  $3.2\times10^{-4}$. Follow-up diagnostics attributed this to mixed-precision
  sensitivity of the finite Newton--Schulz map to tiny perturbations of its
  input, not to a single faulty kernel, and found that FP32 and FP64 controls
  pass while one-ULP perturbations of the inner factor reproduce the failure.
\item An outer-diagonal wave failed a checkpoint-continuation gate: training
  continuously to update 24 and resuming from update 22 produced weight tensors
  differing beyond $\text{rtol}=10^{-5}$, $\text{atol}=10^{-6}$, despite
  matching logged losses. Appendix~\ref{app:numerics} reports what we learned
  chasing this.
\item A production launch died before its first optimizer update on a process
  ownership check, and was re-run as the wave reported in
  Section~\ref{sec:tuning}. The recovery is not an independent replication of
  anything.
\end{itemize}

\subsection{Reproducing the tables and figures}

Every number and every figure in this paper is produced from the per-update
JSONL histories by three scripts, with no manual transcription:
\texttt{scripts/analyze.py} (full-budget runs, the main and window tables, the
cost--quality plane, the paired-delta and freshness series),
\texttt{scripts/analyze\_lr.py} (the 28 screening trials and the sweep figure),
and \texttt{scripts/check\_duplicates.py} (the distinct-trajectory audit in
Table~\ref{tab:evidence}). \texttt{scripts/derive.py} computes the handful of
cross-study comparisons the prose quotes. Running the four in that order
regenerates \texttt{results/} and \texttt{figs/} completely; the manuscript
reads only generated files for run-derived quantities.
